\section{Effective Cartier divisors and invertible sheaves}
\label{section-effective-Cartier-invertible}

\noindent
Since an effective Cartier divisor has an invertible ideal sheaf
(Definition \ref{definition-effective-Cartier-divisor}) the
following definition makes sense.

\begin{definition}
\label{definition-invertible-sheaf-effective-Cartier-divisor}
Let $S$ be a scheme. Let $D \subset S$ be an effective Cartier divisor
with ideal sheaf $\mathcal{I}_D$.
\begin{enumerate}
\item The {\it invertible sheaf $\mathcal{O}_S(D)$ associated to $D$}
is defined by
$$
\mathcal{O}_S(D) =
\SheafHom_{\mathcal{O}_S}(\mathcal{I}_D, \mathcal{O}_S) =
\mathcal{I}_D^{\otimes -1}.
$$
\item The {\it canonical section}, usually denoted $1$ or $1_D$, is the
global section of $\mathcal{O}_S(D)$ corresponding to
the inclusion mapping $\mathcal{I}_D \to \mathcal{O}_S$.
\item We write
$\mathcal{O}_S(-D) = \mathcal{O}_S(D)^{\otimes -1} = \mathcal{I}_D$.
\item Given a second effective Cartier divisor $D' \subset S$ we define
$\mathcal{O}_S(D - D') =
\mathcal{O}_S(D) \otimes_{\mathcal{O}_S} \mathcal{O}_S(-D')$.
\end{enumerate}
\end{definition}

\noindent
Some comments. We will see below that the assignment
$D \mapsto \mathcal{O}_S(D)$ turns addition of effective Cartier
divisors (Definition \ref{definition-sum-effective-Cartier-divisors})
into addition in the Picard group of $S$
(Lemma \ref{lemma-invertible-sheaf-sum-effective-Cartier-divisors}).
However, the expression $D - D'$ in the definition above does not
have any geometric meaning. More precisely, we can think of the
set of effective Cartier divisors on $S$ as a commutative monoid
$\text{EffCart}(S)$ whose zero element is the empty effective Cartier divisor.
Then the assignment $(D, D') \mapsto \mathcal{O}_S(D - D')$ defines
a group homomorphism
$$
\text{EffCart}(S)^{gp} \longrightarrow \Pic(S)
$$
where the left hand side is the group completion of
$\text{EffCart}(S)$. In other words, when we write $\mathcal{O}_S(D - D')$
we may think of $D - D'$ as an element of $\text{EffCart}(S)^{gp}$.

\begin{lemma}
\label{lemma-conormal-effective-Cartier-divisor}
Let $S$ be a scheme and let $D \subset S$ be an effective Cartier divisor.
Then the conormal sheaf is $\mathcal{C}_{D/S} = \mathcal{I}_D|_D =
\mathcal{O}_S(-D)|_D$ and the normal sheaf is
$\mathcal{N}_{D/S} = \mathcal{O}_S(D)|_D$.
\end{lemma}

\begin{proof}
This follows from Morphisms, Lemma \ref{morphisms-lemma-affine-conormal}.
\end{proof}

\begin{lemma}
\label{lemma-ses-add-divisor}
Let $X$ be a scheme. Let $D, D' \subset X$ be effective Cartier divisors
with $D \subset D'$ and let $C \subset X$ be an effective Cartier divisor
such that $D' = D + C$ (which exists by
Lemma \ref{lemma-difference-effective-Cartier-divisors}).
Then there is a short exact sequence
$$
0 \to \mathcal{O}_X(-D)|_C \to \mathcal{O}_{D'} \to \mathcal{O}_D \to 0
$$
of $\mathcal{O}_X$-modules.
\end{lemma}

\begin{proof}
In the statement of the lemma and in the proof we use the equivalence of
Morphisms, Lemma \ref{morphisms-lemma-i-star-equivalence} to think of
quasi-coherent modules on closed subschemes of $X$
as quasi-coherent modules on $X$. Let $\mathcal{I}$ be the ideal
sheaf of $D$ in $D'$ and let $\mathcal{I}_D = \mathcal{O}_X(-D)$ be
the ideal sheaf of $D$ in $X$. Then there is a canonical surjection
$\mathcal{I}_D \to \mathcal{I}$.
Thus it suffices to show the kernel of this map is equal to
$\mathcal{I}_C\mathcal{I}_D$ (with obvious notation).
This can be checked locally, hence we may assume
$X = \Spec(A)$, $D = \Spec(A/fA)$ and $C = \Spec(A/gA)$ where
$f, g \in A$ are nonzerodivisors
(Lemma \ref{lemma-characterize-effective-Cartier-divisor}).
Then $\mathcal{I}_D$ corresponds to $fA$, the nonzerodivisor $fg \in A$
generates the ideal of $D'$, and $\mathcal{I}$ corresponds to $I = fA/fgA$.
Thus the result is clear.
\end{proof}

\begin{lemma}
\label{lemma-invertible-sheaf-sum-effective-Cartier-divisors}
Let $S$ be a scheme.
Let $D_1$, $D_2$ be effective Cartier divisors on $S$.
Let $D = D_1 + D_2$.
Then there is a unique isomorphism
$$
\mathcal{O}_S(D_1) \otimes_{\mathcal{O}_S} \mathcal{O}_S(D_2)
\longrightarrow
\mathcal{O}_S(D)
$$
which maps $1_{D_1} \otimes 1_{D_2}$ to $1_D$.
\end{lemma}

\begin{proof}
Omitted.
\end{proof}

\begin{lemma}
\label{lemma-pullback-effective-Cartier-divisors}
Let $f : S' \to S$ be a morphism of schemes.
Let $D$ be an effective Cartier divisor on $S$.
If the pullback of $D$ is defined then
$f^*\mathcal{O}_S(D) = \mathcal{O}_{S'}(f^*D)$
and the canonical section $1_D$ pulls back to
the canonical section $1_{f^*D}$.
\end{lemma}

\begin{proof}
Omitted.
\end{proof}

\begin{definition}
\label{definition-regular-section}
Let $(X, \mathcal{O}_X)$ be a locally ringed space.
Let $\mathcal{L}$ be an invertible sheaf on $X$.
A global section $s \in \Gamma(X, \mathcal{L})$ is called a
{\it regular section} if the map $\mathcal{O}_X \to \mathcal{L}$,
$f \mapsto fs$ is injective.
\end{definition}

\begin{lemma}
\label{lemma-regular-section-structure-sheaf}
Let $X$ be a locally ringed space. Let $f \in \Gamma(X, \mathcal{O}_X)$.
The following are equivalent:
\begin{enumerate}
\item $f$ is a regular section, and
\item for any $x \in X$ the image $f \in \mathcal{O}_{X, x}$
is a nonzerodivisor.
\end{enumerate}
If $X$ is a scheme these are also equivalent to
\begin{enumerate}
\item[(3)] for any affine open $\Spec(A) = U \subset X$
the image $f \in A$ is a nonzerodivisor,
\item[(4)] there exists an affine open covering
$X = \bigcup \Spec(A_i)$ such that
the image of $f$ in $A_i$ is a nonzerodivisor for all $i$.
\end{enumerate}
\end{lemma}

\begin{proof}
Omitted.
\end{proof}

\noindent
Note that a global section $s$ of an invertible $\mathcal{O}_X$-module
$\mathcal{L}$ may be seen as an $\mathcal{O}_X$-module map
$s : \mathcal{O}_X \to \mathcal{L}$. Its dual is therefore a
map $s : \mathcal{L}^{\otimes -1} \to \mathcal{O}_X$.
(See Modules, Definition \ref{modules-definition-powers}
for the definition of the dual invertible sheaf.)

\begin{definition}
\label{definition-zero-scheme-s}
Let $X$ be a scheme. Let $\mathcal{L}$ be an invertible sheaf.
Let $s \in \Gamma(X, \mathcal{L})$ be a global section.
The {\it zero scheme} of $s$ is the closed subscheme $Z(s) \subset X$
defined by the quasi-coherent sheaf of ideals
$\mathcal{I} \subset \mathcal{O}_X$ which is the image of the
map $s : \mathcal{L}^{\otimes -1} \to \mathcal{O}_X$.
\end{definition}

\begin{lemma}
\label{lemma-zero-scheme}
Let $X$ be a scheme.
Let $\mathcal{L}$ be an invertible sheaf.
Let $s \in \Gamma(X, \mathcal{L})$.
\begin{enumerate}
\item Consider closed immersions $i : Z \to X$ such that
$i^*s \in \Gamma(Z, i^*\mathcal{L})$ is zero
ordered by inclusion. The zero scheme $Z(s)$ is the
maximal element of this ordered set.
\item For any morphism of schemes $f : Y \to X$ we have
$f^*s = 0$ in $\Gamma(Y, f^*\mathcal{L})$ if and only if
$f$ factors through $Z(s)$.
\item The zero scheme $Z(s)$ is a locally principal closed subscheme.
\item The zero scheme $Z(s)$ is an effective Cartier divisor
if and only if $s$ is a regular section of $\mathcal{L}$.
\end{enumerate}
\end{lemma}

\begin{proof}
Omitted.
\end{proof}

\begin{lemma}
\label{lemma-characterize-OD}
\begin{slogan}
Effective Cartier divisors on a scheme are the same as invertible sheaves
with fixed regular global section.
\end{slogan}
Let $X$ be a scheme.
\begin{enumerate}
\item If $D \subset X$ is an effective Cartier divisor, then
the canonical section $1_D$ of $\mathcal{O}_X(D)$ is regular.
\item Conversely, if $s$ is a regular section of the invertible
sheaf $\mathcal{L}$, then there exists a unique effective
Cartier divisor $D = Z(s) \subset X$ and a unique isomorphism
$\mathcal{O}_X(D) \to \mathcal{L}$ which maps $1_D$ to $s$.
\end{enumerate}
The constructions
$D \mapsto (\mathcal{O}_X(D), 1_D)$ and $(\mathcal{L}, s) \mapsto Z(s)$
give mutually inverse maps
$$
\left\{
\begin{matrix}
\text{effective Cartier divisors on }X
\end{matrix}
\right\}
\leftrightarrow
\left\{
\begin{matrix}
\text{isomorphism classes of pairs }(\mathcal{L}, s)\\
\text{consisting of an invertible }
\mathcal{O}_X\text{-module}\\
\mathcal{L}\text{ and a regular global section }s
\end{matrix}
\right\}
$$
\end{lemma}

\begin{proof}
Omitted.
\end{proof}

\begin{remark}
\label{remark-ses-regular-section}
Let $X$ be a scheme, $\mathcal{L}$ an invertible $\mathcal{O}_X$-module,
and $s$ a regular section of $\mathcal{L}$. Then the zero scheme
$D = Z(s)$ is an effective Cartier divisor on $X$ and there are
short exact sequences
$$
0 \to \mathcal{O}_X \to \mathcal{L} \to i_*(\mathcal{L}|_D) \to 0
\quad\text{and}\quad
0 \to \mathcal{L}^{\otimes -1} \to \mathcal{O}_X \to i_*\mathcal{O}_D \to 0.
$$
Given an effective Cartier divisor $D \subset X$ using
Lemmas \ref{lemma-characterize-OD} and
\ref{lemma-conormal-effective-Cartier-divisor}
we get
$$
0 \to \mathcal{O}_X \to \mathcal{O}_X(D) \to i_*(\mathcal{N}_{D/X}) \to 0
\quad\text{and}\quad
0 \to \mathcal{O}_X(-D) \to \mathcal{O}_X \to i_*(\mathcal{O}_D) \to 0
$$
\end{remark}






